\documentclass[11pt,letterpaper]{article}
\usepackage[margin=0.88in]{geometry}
\usepackage[T1]{fontenc}
\usepackage{lmodern}
\usepackage{amsmath,amssymb,amsthm,mathtools,mathrsfs}
\usepackage{booktabs,longtable,array,microtype,enumitem,xcolor,xurl}
\usepackage[colorlinks=true,linkcolor=blue!45!black,urlcolor=blue!45!black,citecolor=blue!45!black]{hyperref}
\usepackage{fancyhdr}
\setlength{\headheight}{14pt}
\pagestyle{fancy}\fancyhf{}
\lhead{\small Parity and Pell witnesses}
\rhead{\small ProveIt research contribution}
\cfoot{\thepage}
\setlength{\parskip}{0.25em}
\setlength{\emergencystretch}{2em}
\setlist{itemsep=0.2em,topsep=0.35em}
\newtheorem{theorem}{Theorem}[section]
\newtheorem{lemma}[theorem]{Lemma}
\newtheorem{proposition}[theorem]{Proposition}
\newtheorem{corollary}[theorem]{Corollary}
\theoremstyle{remark}\newtheorem{remark}[theorem]{Remark}
\newcommand{\Z}{\mathbb Z}\newcommand{\N}{\mathbb N}\newcommand{\Q}{\mathbb Q}
\newcommand{\Con}{\operatorname{Con}}\newcommand{\Refut}{\operatorname{Ref}}
\newcommand{\Pos}{\operatorname{Pos}}\newcommand{\ZFC}{\operatorname{ZFC}}
\newcommand{\degw}{\operatorname{wdeg}}
\newcommand{\file}[1]{\texttt{\detokenize{#1}}}
\title{\textbf{Parity and Pell Witnesses}\\[0.4em]
\large Smaller explicit Diophantine equations for syntactic consistency\\
and a reduced direct proof-calculus quartic}
\author{AI-assisted research prepared for Vladimir Reshetnikov\\
\small ProveIt project}
\date{23 October 2026, Pacific time}
\begin{document}
\maketitle
\begin{abstract}
We give new fixed-coefficient Diophantine constructions whose absence of
integer zeros is provably equivalent in ZFC to syntactic $\Con(\ZFC)$,
by composition with explicitly identified inherited calibration theorems.
The main mechanism uses an odd affine coordinate to combine Pell witnesses.
It excludes cancellation branches that a bare radical norm would otherwise
admit. A two-root version yields $16$ unrestricted integers, exact degree
$112$, and $119$ binary arithmetic operations after a modulus-four
refinement. A three-root version yields $15$ unrestricted integers,
degree $310$, and $136$ operations; a polynomial translation gives
degree $306$ at $139$ operations. A separate sign-preserving
reduction gives $20$ unrestricted integers at degree $18$.

We also reduce the completed direct C222 proof-calculus construction from
$27{,}916$ to $26{,}188$ unrestricted integer variables while retaining
degree four. Its complete $724{,}768$-term polynomial is supplied.
This route uses neither K215 simulation nor a Jones universal polynomial.
The results are explicit improvements and tradeoffs against a pinned
repository and its incoming archives, not a proof of global minimality
or historical priority. The archive contains literal equations, complete
coefficient bytes, mathematical proofs, independent exact checks, and
reproducible builders. No new end-to-end proof-assistant certificate or
exported object-language ZFC derivation is claimed.
\end{abstract}
\tableofcontents
\newpage

\section{What is established}

\subsection{The target and the comparison}
Fix the recursive syntactic calculus for ZFC used in ProveIt. Write
$\Refut_{\ZFC}(p)$ for its finite-refutation predicate and
\[
 \Con(\ZFC)\equiv\neg\exists p\in\N\ \Refut_{\ZFC}(p).
\]
The target is an explicit polynomial $P$ and an ordinary mathematical
proof, formalizable in ZFC, of
\begin{equation}\label{eq:target}
 \ZFC\vdash
 \left(\neg\exists x\in\Z^n\ P(x)=0\right)
 \longleftrightarrow\Con(\ZFC).
\end{equation}
All coordinates in the small equations below range over \emph{all integers}.
A polynomial in fifteen positive integers and one in fifteen unrestricted
integers make different complexity claims.

Number of unknowns, total degree, evaluator size, coefficient size, and
proof complexity need not improve together. Every binary addition,
subtraction, and multiplication costs one operation here, including squares
and scalar multiplication. Reused registers and fixed integer leaves are
free. This circuit convention is not bit complexity.

The comparison snapshot is ProveIt commit
\file{43ac596e36c105a0e8b07c2d1f5ae200cae9c056}. Initial authenticated
plugin reads used the earlier
\file{bf2f392d73eb2ed74bd51cfdb509ecb704d04eb3}. The later snapshot includes
the direct C222 round-twelve archive. The source ledger records every
incoming ZIP hash. We inspected those archives as well as the collection
README, which does not describe all the newer results.

\begin{table}[ht]
\centering\small
\begin{tabular}{@{}p{2.58in}rrr@{}}
\toprule
Construction, all over $\Z$ & Unknowns & Degree & Operations\\
\midrule
Incoming saturated unit-product &16&184&135\\
Incoming saturated sum of squares &16&192&134\\
\textbf{New two-root norm, mod four} &\textbf{16}&\textbf{112}&\textbf{119}\\
\textbf{New three-root norm, mod four} &\textbf{15}&\textbf{310}&\textbf{136}\\
\textbf{New centered norm, mod four} &\textbf{15}&\textbf{306}&\textbf{139}\\
Incoming transport projection &20&20&111\\
\textbf{New rescaled sign guard} &\textbf{20}&\textbf{18}&\textbf{115}\\
\midrule
Incoming direct C222 quartic &27,916&4&---\\
\textbf{New direct C222 quartic} &\textbf{26,188}&\textbf{4}&---\\
\bottomrule
\end{tabular}
\caption{Exact domain-matched comparisons. The sixteen-integer result
improves degree and evaluator size. The other small forms are tradeoffs.
The direct construction is measured by its full sparse polynomial.}
\label{tab:comparison}
\end{table}
The table lists the final modulus-four refinements. We first prove the
slightly simpler odd-coordinate versions, at $120$, $137$, and $140$
operations, and then save one gate in each in Section~\ref{sec:modfour}.
All six literal circuits are retained.

The saturated-certificate archive supplies the calibrated
seventeen-integer parent used below~\cite{sat}. Polynomial shears and
saturation/projection already supply twenty-integer degree-twenty and other
reductions~\cite{shears,projection}. The incoming direct construction already
completes a machine-independent proof-rule route~\cite{round12}; we improve
it rather than claiming that architecture as new.

The repository also has equations in nine positive or ten unrestricted
integer unknowns of much larger degree. General low-variable reductions,
including Sun's recent ten-integer theorem, are separate baselines
\cite{sun,jones}. Our fifteen-integer result is not a fewest-variable record.
It supplies a much lower-degree, short, directly auditable alternative.

\subsection{The new idea and the proof boundary}
The main device is a witness-recovery argument: an odd total coordinate
prevents precisely the cancellation cases that would destroy integrality
in a three-root norm. This is more specific than invoking MRDP.
The resulting objects are fully specified finite polynomials; named
registers are abbreviations, not hidden unknowns.

Three kinds of evidence are distinct:
\begin{enumerate}
\item The new transformations and all exceptional cases are proved below.
\item Exact finite checks bind the formulas to equation files, complete
fixed coefficients, operation counts, and nonzero leading terms.
\item The interpretation as syntactic ZFC refutability uses the named
inherited native/compiler/calculus theorems, or the direct C222 calibration.
\end{enumerate}
The third step is an ordinary mathematical dependency, not a consequence
of tests or hashes. The retained status documents record unfinished
end-to-end native formalization and object-language proof obligations
\cite{status}. Reusing ordinary theorems formalizable in ZFC establishes
the ordinary proof of \eqref{eq:target}; it does not emit a proof file in
an encoded ZFC calculus.

No assertion of $\Con(\ZFC)$ is made. The equations are not claimed to
have been proved unsolvable. The result is the biconditional, whose proof
does not assume consistency.

\section{Fixed coefficients and the calibrated parent}\label{sec:source}

\subsection{Six fixed integers}
The small integer equations use exactly the coherent source-34 tuple:
\[
 \beta=2^{1064921}-1,\quad K_0,\quad M_C,\quad M_F,\quad
 A_{\mathrm c}=2^{1113},\quad B_{\mathrm c}=2^{666104}.
\]
Here $K_0,M_C,M_F$ are fixed numerals, not parameters to choose later.
Their complete values are supplied as minimal unsigned little-endian
byte strings: $b_0,\ldots,b_{\ell-1}$ denotes
$\sum_{j=0}^{\ell-1}b_j256^j$.
The binding records lengths, bit lengths, and SHA-256 values; the bytes,
not the hashes, define the constants. The six raw encodings occupy
$615{,}883$ bytes. The finite recipe, literal machine table, and ordered
native tile data are retained. No coefficient depends on searching for
a refutation.

$M_F$ already includes its inherited shift by $\beta$; it must not be
shifted again. The compact44, compact54, and source-34 ports are not
interchangeable. Every small integer equation here uses the same complete
six-port binding. The inherited source establishes
\[
 \beta\ge17,\quad \beta\equiv1\pmod3,\quad
 0<M_C<\beta<M_F<2\beta,\quad K_0\ge0,\quad
 A_{\mathrm c},B_{\mathrm c}\ge2.
\]
These elementary facts suffice for the new preliminary guards. Calibration
of the actual tuple remains the identified inherited theorem.

\subsection{Shared polynomial registers}
The following coordinates, whenever included in an equation, are
independent integers:
\[
 J,C,z,f,h,v,w,\gamma,T,c,t,p_1,p_2,p_3.
\]
Additional coordinates will be listed for each equation. Define
\begin{align}
 J_0&=3J+1,&q&=\beta J_0+1,&W&=4(3w+1),\notag\\
 X&=qW,&Y&=(2v+1)q^3,&E&=XY,\notag\\
 a&=(X+1)Y,&H&=4a+3,&\Delta&=a^2+H,\label{eq:shared1}\\
 F&=(K_0+W)C-(q-1)z,&\alpha&=q-F-C,\notag\\
 R&=(q-F)(q^3-q)-C(q^2-1)+(M_C+qM_F)J_0,\label{eq:R}\\
 k&=R+1+hE,&\eta&=c-Yk,&\zeta&=k-\eta,\notag\\
 m&=XY^2,&B_0&=mk,&S&=k(2m+1),\label{eq:shared2}\\
 V&=c(Tf-1)-Rf^2,&s&=X+\gamma H.\label{eq:V}
\end{align}
The quotient-compressed $V$ is $cTf-c-Rf^2$, not
$f^2(cT-R)-c$; confusing them would change the polynomial.

Two retained residuals are
\begin{align}
 \mathcal M&=2acs+s^2-Hc^2-1,\label{eq:main}\\
 d_{\mathrm g}&=(A_{\mathrm c}F)(B_{\mathrm c}C)
       -(q-A_{\mathrm c}F-B_{\mathrm c}C)^2,\notag\\
 \Pos&=1+p_1^2+p_2^2+p_3(p_3+1),\notag\\
 \mathcal G&=d_{\mathrm g}\eta\zeta
       -(t^2-d_{\mathrm g}-\eta\zeta+1)^2-\Pos.\label{eq:guard}
\end{align}
Thus $\mathcal M=(ac+s)^2-\Delta c^2-1$.
Every name in these displays is a polynomial abbreviation.

\subsection{The parent and its logical interface}
The inherited parent adds coordinates $r,\tau,\nu$ and sets
$I=(2r+1)c^2$. Its five residuals are
\begin{align}
 \mathcal F&=\tau^2+(2\tau-k)B_0-1,\notag\\
 \mathcal M&,\qquad I^2-\Delta(f^2-1),\notag\\
 \mathcal A&=(V+\nu)^2-I^2\nu(2V+\nu)-1,\notag\\
 \mathcal G&.\label{eq:parent}
\end{align}
Their squared sum is its defining polynomial. Put
\[
 \tau_{\rm old}=\tau+B_0,\qquad y=V+\nu,\qquad
 A_0=1+m(m+1)k^2.
\]
The exact centering identities are
\begin{align}
 \mathcal F&=\tau_{\rm old}^2-A_0,\label{eq:firstcenter}\\
 \mathcal A&=I^2V^2-(I^2-1)y^2-1.\label{eq:auxcenter}
\end{align}

The inherited theorem identifies existence of an integer zero of this
literal parent with finite syntactic ZFC refutability~\cite{sat}.
Its proof restores the signed native/Pell source, including the quotient
witnesses $V=of-c=jc-R$. The required divisibility $c\mid f^2-1$ is
proved by the ordinary strong-rank argument before those omitted quotients
are used. This ordering is essential.

Our reductions keep $\mathcal M$ and $\mathcal G$ and recover exactly the
other three parent rows. Their existence equivalence is to this particular
calibrated parent, not to an unspecified universal polynomial.

\section{Lemmas available before root recovery}\label{sec:lemmas}

\begin{lemma}[Inherited joint positivity]\label{lem:joint}
For integers $x,y$,
\[
 x>0,\ y>0\ \Longleftrightarrow\
 \exists t,p_1,p_2,p_3\in\Z:\
 xy-(t^2-x-y+1)^2=\Pos(p_1,p_2,p_3).
\]
\end{lemma}
\begin{proof}
The right side is positive, hence $xy>0$. If $x=-a,y=-b$ with $a,b>0$,
then $(t^2+a+b+1)^2>ab$, a contradiction.

Conversely put $N=x+y-1\ge1$. Since $xy-N=(x-1)(y-1)\ge0$, choose a
neighboring square $t^2$ with $(t^2-N)^2<N$. If $N$ is a square the
error is zero; otherwise one of the squares adjacent to $N$ has error
at most $\lfloor\sqrt N\rfloor$, whose square is less than $N$.
The resulting positive difference is a value of $\Pos$:
Legendre's three-square theorem represents $4n-3$, which is $1$ or $5$
modulo $8$. Exactly one square is odd. Writing it as $(2p_3+1)^2$
and the other two as $(2p_1)^2,(2p_2)^2$ gives the formula.
\end{proof}

\begin{lemma}[Guard-only nonzero scales]\label{lem:guard}
If $\mathcal G=0$, then
\[
 |k|\ge2,\quad c\ne0,\quad\Delta>1,\quad A_0>0,\quad
 S\ne0,\quad S^2\ge4A_0>A_0.
\]
No first, strong, or auxiliary Pell equation is used here.
\end{lemma}
\begin{proof}
Lemma~\ref{lem:joint} gives $d_{\mathrm g}>0$ and $\eta\zeta>0$.
Thus $\eta,\zeta$ have the same nonzero sign, $|k|\ge2$, and
$|\eta|<|k|$. The integer $J_0=3J+1$ is nonzero; the actual $\beta$
therefore gives $q\ne0$ and $|q|\ge16$.
Also $W=12w+4$ has $|W|\ge4$, so $|X|\ge64$ and $|Y|\ge16^3$.
It follows that
\[
 |c|=|\eta+Yk|>|k|(|Y|-1)>0,\qquad
 |a+2|=|(X+1)Y+2|>2.
\]
Hence $\Delta=(a+2)^2-1>1$.
For every integer $m$, $m(m+1)\ge0$, so $A_0\ge1$.
Finally the exact identity
\begin{equation}\label{eq:Ssep}S^2-4A_0=k^2-4\end{equation}
gives the asserted separation.
\end{proof}

We retain the adjacent-coefficient descent from the saturated source
and its package-69 predecessor~\cite{sat}; it is reproduced here to make
the new root-recovery argument self-contained.
\begin{lemma}[Inherited adjacent-coefficient rational descent]\label{lem:descent}
If $Q>1,V$ are integers and $y\in\Q$ satisfies
$QV^2-(Q-1)y^2=1$, then $y\in\Z$.
\end{lemma}
\begin{proof}
Write $y=a/b$ in lowest terms with $b>0$. Clearing denominators gives
$b^2\mid Q-1$. Take absolute values of $V,y$. The norm excludes $V=0$;
at $V=1$ it gives $y=1$ and $b=1$.

For $V>1$ use the norm-preserving transformation
\[
 V'=(2Q-1)V-2(Q-1)y,\qquad y'=(2Q-1)y-2QV.
\]
The first coordinate is integral because $b^2\mid Q-1$.
The reduced denominator of $y'$ is exactly $b$, because
$2Q-1\equiv1\pmod b$. The norm is preserved by direct expansion.
Since $y>V$, $V'<V$. Also
\[
 Q^2V^2-(Q-1)^2y^2=QV^2+Q-1>0,
\]
so $(Q-1)y<QV$ and $V'>-V$. The transformed norm excludes $V'=0$.
Thus $0<|V'|<V$. Repeated absolute-value descent through positive
integers preserves the denominator until $|V|=1$, where it is one.
\end{proof}
This is a theorem for adjacent coefficients, not for arbitrary rational
points on arbitrary Pell conics. Its unbounded proof cannot be replaced
by a finite rational-root search.

\section{An odd two-root norm: sixteen integers, degree 112}
\label{sec:sixteen}

\begin{theorem}[Odd two-root extraction]\label{thm:two}
Let $A,Q,d,L\in\Z$, with $d\ne0$ and $L$ odd. Put
\[
 B=d(L^2-4A)+Q,\qquad N_2=B^2-4dL^2Q.
\]
Then $N_2=0$ if and only if there are an odd integer $u$ and an integer
$s$ satisfying $du^2=Q$, $s^2=A$, and $L=u+2s$.
\end{theorem}
\begin{proof}
For soundness, $L\ne0$. Set $u=B/(2dL)\in\Q$.
The norm gives $u^2=Q/d$; the definition of $B$ then gives
$(L-u)^2=4A$. Thus $s=(L-u)/2$ is a rational square root of an integer,
so is integral. Indeed, for a reduced fraction its denominator squared
must divide its numerator squared. Then $u=L-2s$ is an odd integer.
Conversely, after $Q=du^2$ the identity
\[
 N_2=d^2\bigl((L-u)^2-4A\bigr)\bigl((L+u)^2-4A\bigr)
\]
proves the claim.
\end{proof}
Parity both prevents a zero extraction denominator and restores the
omitted multiplier in the required odd class.

Use the sixteen coordinates
\[
 J,C,z,f,h,n_0,v,w,\nu,\gamma,T,c,t,p_1,p_2,p_3\in\Z.
\]
In addition to the shared definitions, set
\begin{align}
 n&=2n_0+1,&d&=c^4,&Q&=\Delta(f^2-1),\notag\\
 L&=n+2B_0,&G_0&=n^2+4B_0(n-k)-4,\notag\\
 B&=dG_0+Q,&N_2&=B^2-4dL^2Q,\label{eq:16norm}\\
 \mathcal A_2&=(V+\nu)^2-Q\nu(2V+\nu)-1.\notag
\end{align}
The complete equation is
\begin{equation}\label{eq:P16}
 \boxed{P_{16}=\mathcal M^2+\mathcal G^2+\mathcal A_2^2+N_2^2=0.}
\end{equation}
The identity $G_0=(n+2B_0)^2-4A_0$ cancels the largest first-Pell
terms before evaluation.

\begin{theorem}\label{thm:16}
$P_{16}$ has an integer zero if and only if the parent
\eqref{eq:parent} does. Consequently it satisfies \eqref{eq:target}.
\end{theorem}
\begin{proof}
At a new zero each residual vanishes, and the guard gives $c\ne0$.
Apply Theorem~\ref{thm:two} with $A=A_0,d=c^4$ to obtain
$du^2=Q$, $\tau_{\rm old}^2=A_0$, and $L=u+2\tau_{\rm old}$,
with $u$ odd and $\tau_{\rm old}$ integral. Restore
\[
 r=(u-1)/2,\qquad \tau=\tau_{\rm old}-B_0,\qquad I=uc^2.
\]
The first and strong parent rows vanish. Since $I^2=Q$, the retained
auxiliary row is exactly the old one. The main and guard rows are unchanged.

Conversely take a parent zero and set $n_0=r+\tau$.
Then $L=(2r+1)+2(\tau+B_0)$, so the two-root identity gives $N_2=0$.
All other new rows vanish. Compose with the parent/refutation theorem.
\end{proof}

The literal division-free circuit has $120$ live operations. The core
norm has exact degree $56$, the guard degree $44$, and the other
residuals smaller degree. Hence $P_{16}$ has exact degree $112$.
A nonzero coefficient certificate is
\[
 [c^{16}J^{60}w^{12}v^{16}h^8]P_{16}
   =2^{48}3^{72}\beta^{60}>0.
\]
The verifier checks the exact upper bound and this coefficient after
specializing the actual tuple. Compared with the incoming $16/184/135$
point, degree falls by $72$ and evaluator size by $15$ operations.
This two-root proof does not require rational Pell descent.

\section{Odd three-root recovery}\label{sec:three}

\subsection{The algebraic norm}
For rational $a,b,c$, define
\begin{equation}\label{eq:norm3}
 \mathscr N(L;a,b,c)
 =\bigl((L^2-a-b-c)^2-4(ab+ac+bc)\bigr)^2-64abcL^2.
\end{equation}
Over a field containing their square roots, this is the exact product
\[
 \prod_{\epsilon_1,\epsilon_2,\epsilon_3\in\{-1,1\}}
 (L-\epsilon_1\sqrt a-\epsilon_2\sqrt b-\epsilon_3\sqrt c).
\]
A zero supplies a signed sum of roots, but does not by itself make the
individual roots rational. The next argument removes that obstruction.

\begin{lemma}\label{lem:classes}
If a signed sum of square roots of three positive rationals is a nonzero
rational, either all roots are rational, or exactly two irrational roots
cancel and the remaining root is rational.
\end{lemma}
\begin{proof}
In the finite multiquadratic extension generated by the roots, group the
terms by their square class in $\Q^\times/(\Q^\times)^2$.
Distinct square classes have linearly independent square roots over $\Q$,
as follows from the character basis of the multiquadratic extension.
The trivial square class supplies the rational part, so at least one
term is rational. If any term is irrational, the other irrational term
must lie in the same class and cancel it.
\end{proof}

\begin{theorem}[Odd three-root Pell merger]\label{thm:three}
Let $A,Q,S,V,c,L$ be integers such that
\[
 A>0,\quad c\ne0,\quad S^2>A,\quad Q\ne1,\quad L\text{ is odd}.
\]
Put $d=c^4$ and
\[
 a=4A,\qquad b=Q/d,\qquad c_*=\frac{4S^2(QV^2-1)}{Q-1}.
\]
Then $\mathscr N(L;a,b,c_*)=0$ if and only if there are integers
$u,\tau,y$, with $u$ odd, satisfying
\begin{equation}\label{eq:three-witnesses}
 du^2=Q,\qquad \tau^2=A,\qquad
 QV^2-(Q-1)y^2=1,\qquad L=u+2\tau+2Sy.
\end{equation}
\end{theorem}
\begin{proof}
Suppose the norm vanishes. Its factorization supplies some signed root
sum. Establish the signs before applying a positive-radical argument.

\emph{If $Q<0$,} $a,c_*>0$ but $b<0$. The signed sum has exactly one
nonzero purely imaginary term and two real terms, and cannot equal real $L$.
This is impossible.

\emph{If $Q=0$,} the factorization gives $L=\pm2\sqrt A\pm2S$.
Its rationality makes $\sqrt A$ rational, hence integral. The right side
is even, contrary to odd $L$.
The value $Q=1$ is excluded by hypothesis; hence $Q>1$.
If $V=0$, then $a,b>0$ and $c_*<0$, again giving one nonzero imaginary
term. Therefore $|V|\ge1$.

Now put $D_y=(QV^2-1)/(Q-1)\ge1$. All radicands are positive and
\[
 c_*=4S^2D_y\ge4S^2>4A=a.
\]
In particular the first and third radicals cannot cancel.

By Lemma~\ref{lem:classes}, if not all roots are rational, two irrational
ones must cancel. If the first and second cancel, the remaining rational
root is $2Sy$ with rational $y$ on
$QV^2-(Q-1)y^2=1$. Lemma~\ref{lem:descent} gives $y\in\Z$,
making $L$ even. If the second and third cancel, the remaining root
is $2\tau$ with $\tau^2=A$, so $\tau\in\Z$ and $L$ is even.
The first and third cannot cancel by the strict inequality above.
All possibilities contradict oddness. Thus all three roots are rational.

Choose the signs supplied by the vanishing norm factor. Rational-square
integrality gives $\tau\in\Z$, and Lemma~\ref{lem:descent} gives $y\in\Z$.
For either sign of nonzero $S$, the third signed root may be written
$2Sy$. The remaining root is $u=L-2\tau-2Sy$, an odd integer whose
square is $Q/d$. This proves \eqref{eq:three-witnesses}.

Conversely, those equations supply one zero factor in the displayed
factorization of $\mathscr N$.
\end{proof}

\begin{remark}
The strong radicand $Q/d$ is allowed to coincide with either other
radicand. Such equality is a genuine loophole for a bare norm construction.
Here any cancellation involving it forces the prescribed odd coordinate
to be even. Only the two even-weighted roots require strict separation.
\end{remark}

\section{The fifteen-integer polynomial}\label{sec:fifteen}

\subsection{A complete division-free definition}
The independent coordinates are exactly
\[
 \boxed{J,C,z,f,h,n_0,v,w,\gamma,T,c,t,p_1,p_2,p_3\in\Z.}
\]
Use the shared registers of Section~\ref{sec:source} and define
\begin{align}
 A_0&=1+m(m+1)k^2,&n&=2n_0+1,&d&=c^4,\notag\\
 Q&=\Delta(f^2-1),&\delta&=Q-1,&D&=d\delta,\notag\\
 L&=n+2B_0+2SV,\label{eq:15L}\\
 \overline A&=4A_0D,&\overline B&=Q\delta,&
 \overline C&=4S^2d(QV^2-1),\notag\\
 U&=DL^2-\overline A-\overline B-\overline C,\notag\\
 T_3&=U^2-4(\overline A\overline B+\overline A\overline C+
                  \overline B\overline C),\notag\\
 N_3&=T_3^2-64\overline A\overline B\overline C\,DL^2,\label{eq:15norm}\\
 G&=\mathcal M^2+\mathcal G^2.\notag
\end{align}
The equation is
\begin{equation}\label{eq:P15}
 \boxed{P_{15}=G+N_3+GN_3=0.}
\end{equation}
All products and exponents are fixed. The rational expressions in
Theorem~\ref{thm:three} occur only in its proof; the polynomial has no
division. The exact denominator-clearing identity is
\[
 N_3=D^4\mathscr N
 \left(L;4A_0,\frac Qd,\frac{4S^2(QV^2-1)}{Q-1}\right)
\]
whenever $D\ne0$.

\subsection{The unit product and the complete witness proof}
The finalizer saves degree:
\[
 P_{15}=(1+G)(1+N_3)-1.
\]
Since $G$ is a sum of integer squares, $1+G$ is a positive integer.
If its product with $1+N_3$ equals one, then $1+G=1$ and
$N_3=0$. Conversely $G=N_3=0$ makes the equation vanish.
No assumption that $N_3$ is nonnegative is made.

\begin{theorem}\label{thm:15}
$P_{15}$ has an integer zero if and only if the calibrated parent
\eqref{eq:parent} does. In particular
\[
 \ZFC\vdash
 \left(\neg\exists x\in\Z^{15}:P_{15}(x)=0\right)
 \longleftrightarrow\Con(\ZFC).
\]
\end{theorem}
\begin{proof}
At a new zero, $\mathcal M=\mathcal G=N_3=0$. The guard gives
$c\ne0,\Delta>1,A_0>0,S^2>A_0$.
Also $Q=\Delta(f^2-1)\ne1$, because an integer multiple of
$\Delta>1$ cannot be one. Hence $D=c^4(Q-1)\ne0$.
The coordinate $L$ is odd.

Apply Theorem~\ref{thm:three} to the cleared norm. Recover odd $u$ and
integers $\tau_{\rm old},y$, and restore
\[
 r=(u-1)/2,\qquad \tau=\tau_{\rm old}-B_0,\qquad \nu=y-V.
\]
With $I=uc^2$, the recovered $I^2=Q$ is the parent strong equation.
Identities \eqref{eq:firstcenter} and \eqref{eq:auxcenter} restore its
first and auxiliary equations. The other two were retained.

Conversely, from a parent zero choose
\begin{equation}\label{eq:15map}n_0=r+\tau+S\nu.\end{equation}
Then, with $u=2r+1$, $\tau_{\rm old}=\tau+B_0$ and $y=\nu+V$,
we have $L=u+2\tau_{\rm old}+2Sy$. One norm factor vanishes, so
$N_3=0$. The retained rows give $G=0$ and hence $P_{15}=0$.
The inherited parent/refutation theorem gives the claimed ZFC equivalence.
\end{proof}

The soundness proof uses rational expressions only as proof tools and
then proves the recovered values integral. The reverse map
\eqref{eq:15map} is polynomial in the old integer coordinates.
Existential equivalence is the relevant statement; no bijection between
the complete witness sets is needed.

\subsection{Exact degree and operation count}
The literal circuit has $137$ live binary operations.
Temporarily assign weights
\[
 \degw(m,k,V,Q,d,n)=(10,7,7,14,4,1).
\]
The raw expression suggests degree $118$ for $T_3$, but its leading
terms cancel. Exact expansion with $S=k(2m+1)$ and
$A_0=1+m(m+1)k^2$ gives
\[
 \degw T_3=111,\qquad
 \operatorname{in}(T_3)=-512m^4k^4V^3Qd^2.
\]
The other term of $N_3$ has degree at most $212$. Consequently
\begin{equation}\label{eq:15leader}
 \degw N_3=222,\qquad
 \operatorname{in}(N_3)=2^{18}m^8k^8V^6Q^2d^4.
\end{equation}
This is an exact $3{,}120$-term polynomial computation, not a modular
estimate. Substitution of the source registers retains degree $222$.
The highest-degree part of $G$ has degree $88$ and its product with
\eqref{eq:15leader} is nonzero, giving degree $310$ for $P_{15}$.
One actual coefficient certificate is
\[
 [c^{16}J^{184}w^{28}v^{44}h^{16}z^6f^{16}]P_{15}
 =2^{118}3^{212}\beta^{184}>0.
\]
The exponents sum to $310$. No genericity condition on $K_0,M_C,M_F$
is required.

The sum-of-squares alternative $G+N_3^2=0$ has degree $444$.
The positive-unit finalizer saves $134$ degrees at one additional binary
operation. Its validity is the integer-unit argument above.

\section{A lower-degree polynomial translation}\label{sec:centering}
Replace only the definition of $n$ by
\begin{equation}\label{eq:centered-n}n=2(n_0-hYR)+1.\end{equation}
This is an invertible polynomial change in the integer coordinate $n_0$:
$h,Y,R$ do not involve that coordinate. Oddness is unchanged, and the
completeness map becomes
$n_0=r+\tau+S\nu+hYR$.

The motivation is the leading correction in the near-root branch:
$SV/Q$ has polynomial leading part $-2hYR$ after the source substitutions.
The observation suggests a translation; exact expansion proves its effect.
The certificate uses the authentic expressions
\[
 m=XY^2,\quad k=R+1+hXY,\quad V=c(Tf-1)-Rf^2,
\]
\[
 Q=((X+1)^2Y^2+4(X+1)Y+3)(f^2-1).
\]
Give $X,Y,R$ weights $2,4,5$, and $h,f,c,T,n_0$ weight one.
The inner norm expands to $61{,}891$ terms, of exact weighted degree
$109$, with leading form
\[
 512X^9Y^{13}R^3h^3f^6c^8
       \bigl[Rf^2+hY(X-2f^2)\bigr].
\]
The subtracted triple-radical tail has degree at most $212$.
Thus the norm has degree $218$, and the final polynomial has
degree $88+218=306$. On the actual source substitution a certificate is
\[
 [J^{184}w^{26}v^{42}z^8h^{14}f^{16}c^{16}]P_{15,\mathrm c}
   =2^{112}3^{210}\beta^{184}>0.
\]
The translated circuit has $140$ operations. Evaluating
$2(n_0-hYR)+1$ costs three more operations than $2n_0+1$.
Both fifteen-variable equations are retained: degree $310$ at $137$
operations and degree $306$ at $140$ operations.

\section{The final modulus-four refinement}\label{sec:modfour}

The authentic source has a useful residue property without any Pell
equation: $W=4(3w+1)$ makes $X$, and hence $m=XY^2$, divisible by four.
Consequently
\[
 B_0=mk\equiv0\pmod4,\qquad A_0=1+m(m+1)k^2\equiv1\pmod4.
\]
Any integer square root of $A_0$ is therefore odd. Prescribing a sum
congruent to two modulo four can replace the factors of two used in the
preceding first-root encodings.

\subsection{Sixteen integers at 119 operations}
Replace \eqref{eq:16norm} by
\begin{align*}
 n&=4n_0+2,&L&=n+B_0,&d&=c^4,&Q&=\Delta(f^2-1),\\
 G_4&=n^2+B_0(2n-k)-1=L^2-A_0,\\
 B_4&=dG_4+Q,&N_{2,4}&=B_4^2-4dL^2Q.
\end{align*}
Keep $\mathcal M,\mathcal G,\mathcal A_2$ and use their squared sum
with $N_{2,4}^2$. At a zero, $L\equiv2\pmod4$, so $L\ne0$.
Extract $u=B_4/(2dL)$ and $\tau_{\rm old}=L-u$.
The norm identity gives $du^2=Q$ and $\tau_{\rm old}^2=A_0$.
Thus $\tau_{\rm old}$ is integral and odd, and $u=L-\tau_{\rm old}$
is also an odd integer. The same parent restoration applies.

For completeness, fix the parent strong multiplier $u=2r+1$.
Choose the sign of the private first root so that
$\tau_{\rm old}\equiv u\pmod4$. This is always possible for an odd root
and changes no other parent equation. Then
$u+\tau_{\rm old}\equiv2\pmod4$, and
\[
 n_0=(u+\tau_{\rm old}-B_0-2)/4
\]
is integral. It supplies the new zero. This is an existence-preserving
sign choice, not a claim that every old value of the centered first-root
coordinate is kept.

The circuit uses $119$ operations. Its exact degree remains $112$.
The coefficient of $c^{16}J^{60}w^{12}v^{16}h^8$ is now
$2^{40}3^{72}\beta^{60}$.

\subsection{Fifteen integers at 136 or 139 operations}
In the three-root definitions use
\[
 n=4n_0+2,\qquad L=n+B_0+4SV,
\]
and replace only the first and third scaled radicands by
\[
 \overline A=A_0D,\qquad
 \overline C=16S^2d(QV^2-1).
\]
Keep $\overline B=Q(Q-1)$, the same norm expression
\eqref{eq:15norm}, and the unit finalizer $G+N_3+GN_3$.
The rational radicands are now
\[
 A_0,\qquad Q/d,\qquad
 16S^2(QV^2-1)/(Q-1),
\]
and $L\equiv2\pmod4$.

The negative-radicand exclusions are unchanged. At $Q=0$, the signed
sum is $\pm\sqrt{A_0}\pm4S$, which is odd if rational, and cannot equal
even $L$. In the positive case, first/third cancellation is again excluded
by separation. First/second cancellation would leave $4Sy$, divisible by
four after rational Pell descent; second/third cancellation would leave
the odd integer first root. Neither can equal $L\equiv2\pmod4$.
Thus all roots are rational. The recovered first root is odd,
$y$ is integral, and
$u=L-\tau_{\rm old}-4Sy$ is an odd integer.
The parent restoration is complete.

For completeness, again choose the sign of the private first root so
$u+\tau_{\rm old}\equiv2\pmod4$, and set
\[
 n_0=(u+\tau_{\rm old}-B_0-2)/4+S\nu.
\]
This proves equivalence without an extra condition on a canonical
computation witness.

The uncentered polynomial has $136$ operations and degree $310$.
Replacing $n$ by $4(n_0-hYR)+2$ gives $139$ operations and degree $306$;
the completeness map simply adds $hYR$. The corresponding coefficient
certificates are
\begin{align*}
 [c^{16}J^{184}w^{28}v^{44}h^{16}z^6f^{16}]P_{15,4}
   &=2^{122}3^{212}\beta^{184},\\
 [J^{184}w^{26}v^{42}z^8h^{14}f^{16}c^{16}]P_{15,4,\mathrm c}
   &=2^{116}3^{210}\beta^{184}.
\end{align*}
The independent centered expansion has $61{,}557$ terms and inner
degree $109$. Its leading form is four times the one in
Section~\ref{sec:centering}; the norm degree is $218$.
These are the final small-equation evaluators recommended in
Table~\ref{tab:comparison}.

\section{Rescaling a sign condition: twenty integers, degree 18}
\label{sec:twenty}
The parent of this independent construction is the active centered joint22
equation~\cite{joint22}. A new coordinate $b=a-E$ allows the first two
source rows to restore
$E=Xb$ and $a=(X+1)b$, retaining
\[
 b=(2v+1)q^3.
\]
Thus $b\ne0$ from the actual $\beta>1$, the form of $q$, and the odd
factor. The monic transport pivot replaces $(F,\alpha)$ by $C$ and
$F=(K_0+W)C-z(q-1)$.

A naive substitution raises the joint guard's degree unnecessarily.
Its old positive argument is $\eta(k-\eta)$, while a retained row gives
$c-\eta=kb$. Set
\[
 u=b\eta,\qquad t'=u(c-\eta-u).
\]
On that retained row,
\begin{equation}\label{eq:signscale}t'=b^2\eta(k-\eta).\end{equation}
Since $b\ne0$, positivity of the new and old arguments is equivalent.
Use Lemma~\ref{lem:joint} on $(d_{\mathrm g},t')$.
This preserves existence of private guard witnesses, not their old values:
the nearest-square lemma regenerates the four private coordinates in each
direction. It is not an equality of the two guard polynomials.

Use the strong row to replace $\Delta(f^2-1)$ by $I^2$ in the
centered auxiliary row before eliminating the strong ordinate. The row is
\[
 (fo+y)^2-I^2(c+y)(2fo+y-c)-1,\qquad I=(2r+1)c^2.
\]
The eight residual degrees are $4,5,8,8,5,8,9,8$.
The unique degree-nine component is $-8r^2c^4fo(c+y)$.
The full sum of squares therefore has exact degree eighteen and
highest-degree form
\[
 64r^4c^8f^2o^2(c+y)^2.
\]
This certificate is independent of all fixed ports.
The circuit has $115$ operations: $57$ multiplications, $30$ additions,
and $28$ subtractions. Compared with the incoming $20/20/111$ equation,
it trades four gates for a ten-percent degree decrease.
The packet contains complete residual definitions, exact parent restorations, and
the proof of every arithmetic row and \eqref{eq:signscale}.

\section{A smaller completed direct proof-rule quartic}
\label{sec:direct-normalized-power}

The incoming direct C222 construction has a particularly useful feature:
its coefficients are ordinary printed integers, and the entire final
polynomial is explicitly materialized. It uses a finite proof-flow
certificate for twelve numerical Horn rules, followed by a CRT bounded
family and finite arithmetic representations. No Turing-machine simulation
or Jones universal polynomial occurs in this route. The incoming
fixed-base affine construction has $16036$ natural unknowns, or $27916$
unrestricted integer unknowns after its nonnegativity conversion.

We reduce these counts to $14956$ and $26188$, respectively, without
altering the twelve rules, the proof-flow encoding, or the CRT architecture.
The change concerns a redundant truncated-subtraction certificate inside
each of the $108$ calls to the power representation.

\subsection{The rank already present in the power certificate}

All variables in this subsection range over $\mathbb N=\{0,1,2,\ldots\}$.
Write the Pell residual as
\[
 \mathcal N(A;X,Y)=X^2+Y^2-A^2Y^2-1.
\]
The active finite branch of the inherited power representation contains
three facts:
\begin{equation}\label{eq:norm-rank-hypotheses}
 a>1,\qquad n\leq w,\qquad
 \mathcal N(w+1;a,wz)=0.
\end{equation}
The first follows from the outer positive-parameter guard in the Pell
representation, including when that representation uses its zero-index
branch. The second and third are explicit residuals in the power branch.

\paragraph{Rank lemma.}
Conditions~\eqref{eq:norm-rank-hypotheses} imply
\[
 w\geq1,\qquad z\geq1,\qquad a>w\geq n.
\]
Indeed, the norm equation is precisely
\[
 a^2=1+w^3(w+2)z^2.
\]
If $w=0$ or $z=0$, it gives $a^2=1$, contradicting $a>1$.
Consequently $w,z\geq1$, and
\[
 a^2=1+w^2\bigl(w(w+2)z^2\bigr)>w^2.
\]
The claimed inequality follows by monotonicity of squaring on naturals.
This is an unbounded arithmetic proof, not an inference from a finite search.

\subsection{Eliminating two semantic witnesses}

The original finite branch introduces a natural $d$ and a natural slack $s$
through the nonnegative disjunction polynomial
\begin{equation}\label{eq:old-truncated-subtraction}
 (a-d-n)^2\bigl((n-a-s)^2+d^2\bigr)=0.
\end{equation}
It expresses $d=\max(a-n,0)$. Under the rank lemma, the second factor
cannot vanish, since $a>n$ makes $n-a-s\ne0$. Thus $d=a-n$.
The same branch contains the congruence, written with natural quotient
witnesses $q_0,q_1$,
\[
 x+tq_0=yd+m+tq_1.
\]
Substitute $d=a-n$ and move $yn$ to the left. The resulting equation is
\begin{equation}\label{eq:translated-power-congruence}
 x+yn+tq_0=ya+m+tq_1.
\end{equation}
Both sides are natural arithmetic expressions. The new compiler deletes
$d$, $s$, and~\eqref{eq:old-truncated-subtraction}, and replaces the old
congruence by~\eqref{eq:translated-power-congruence}. Every other byte of
the original compiler is retained; the delivered delta checker verifies this.

\paragraph{Soundness and completeness, including inactive branches.}
When the finite branch is selected, its residuals form a sum of
nonnegative expressions, so each vanishes. The rank lemma applies to the
new branch as well as the old one. An old witness therefore restricts to a
new witness. Conversely, a new finite-branch witness extends to an old one
by setting $d=a-n$ and $s=0$. The first factor in
\eqref{eq:old-truncated-subtraction} is then zero, and the two congruences
are the same integer-polynomial equation.

The inherited power representation separately covers exponent zero and
zero base with positive exponent. Neither branch is changed. When one of
these branches is selected, the finite expression is underneath a zero
factor in the disjunction, so it imposes no condition. The erased witnesses
can then be restored as $d=s=0$. All deleted or changed natural arithmetic
wires can be recomputed from their inputs. There is no requirement that an
unselected subtraction or Pell branch vanish. This observation is essential
to the two-way existential equivalence of the complete power graph.

\subsection{Composition into one literal polynomial}

Replacing an existential arithmetic graph by an equivalent graph preserves
every enclosing binomial, factorial and bounded-CRT representation. The
unchanged global construction uses $108$ power graphs, $33$ binomial graphs,
$3$ factorial graphs and $72$ positive-divisor Euclidean-division graphs.
All twelve local Horn branches and both global instantiations of the local
row remain unchanged.

Let $\rho_1,\ldots,\rho_{11502}$ be the explicit quadratic residuals in
\texttt{residuals\_affine.json}, in $14956$ natural variables. The natural
equation is
\[
 P_{\mathbb N}(X)=\sum_{j=1}^{11502}\rho_j(X)^2=0.
\]
The residuals are literal finite integer polynomials, supplied in sparse
form, and their complete expansion is \texttt{equation\_affine.json}.
There are $3744$ semantic inputs and $11212$ retained nonlinear wires.

To obtain the unrestricted-integer equation, introduce three fresh integer
variables $u_i,v_i,w_i$ for each semantic input $X_i$, $0\leq i<3744$, and set
\[
 P_{\mathbb Z}
 =P_{\mathbb N}
  +\sum_{i=0}^{3743}
    \bigl(X_i-u_i^2-v_i^2-w_i(w_i+1)\bigr)^2.
\]
The polynomial has $14956+3\cdot3744=26188$ integer variables and total
degree exactly four. Every natural is of the form
$u^2+v^2+w(w+1)$ with integer $u,v,w$: apply the three-square theorem to
$4X+1$, whose residue modulo eight is one or five, and separate the single
odd square from the two even squares. Conversely, the displayed expression
is nonnegative for every integer triple. All other wires are forced to be
nonnegative by the unchanged topological equations and nonnegative affine
forms, so they need no additional guards.

The complete integer expansion is \texttt{equation\_integer.json}; it has
$724768$ nonzero monomials and maximum absolute coefficient
$320307372009864$. It contains no variable-exponent, factorial, binomial,
bounded-quantifier or symbolic-coefficient operation.

\subsection{Logical calibration and evidence boundary}

The unchanged proof-flow certificate is equivalent to a finite derivation of
the distinguished target in C222. The pinned logical interfaces are
\texttt{KEndpoint.c222\_iff\_zf\_refutation} and
\texttt{ConstructibleChoice.zf\_refutation\_iff\_zfc\_refutation}.
The fixed Horn coefficients are independently matched to
\texttt{Code254.lean} and \texttt{KCode222.lean}; the names of those files
do not introduce a Turing-machine simulation into this construction.
The inherited uniform CRT/Pell proof, the normalization proved above, and
the finite circuit lifting yield ordinary mathematical proofs formalizable
in ZFC of
\[
 \operatorname{ZFC}\vdash
 \bigl(\neg\exists X\in\mathbb N^{14956}:P_{\mathbb N}(X)=0\bigr)
 \;\longleftrightarrow\;\operatorname{Con}(\operatorname{ZFC}),
\]
and the corresponding statement for $P_{\mathbb Z}$ over
$\mathbb Z^{26188}$. No consistency assumption is used in either direction.

The exact finite audits check all arithmetic-interface polynomials,
compositions, Horn branches, source substitutions, gate residuals,
nonnegative affine eliminations and full expanded polynomial identities.
The original arithmetic and logical source files are supplied with immutable
hashes. These checks do not amount to a newly compiled Lean or Rocq proof of
the complete construction. The claimed result is an ordinary proof with
fully materialized and exactly audited polynomial artifacts.

This is a strict reduction of the incoming completed direct construction.
It is not a global fewest-variable result, and it does not dominate the
repository's native low-variable equations. The elementary rank argument is
not claimed to be historically new; its useful application here is the
removal of a redundant branch certificate from every power call.


\section{Why the direct equation needs no machine}
\label{sec:directflow}
The direct component offers a complementary notion of simplicity: its
coefficients are ordinary modest integers, with no native coefficient
payload. Its predecessor already supplies the direct architecture; the new
power normalization reduces the complete measured output.

\subsection{Unordered Horn proof certificates}
Each numerical C222 rule has one conclusion and at most two premises,
padded by truth label $0$. A finite proof tree gives a nonnegative sum of
rule balance vectors equal to $e_1$, counting positive labels and ignoring
label $0$.

Conversely, close the finitely supported rules under derivation from their
zero-premise instances. If target $1$ lay outside the closure, sum the
balance coordinates outside it. Every rule contributes at most zero:
an outside conclusion needs an outside premise, otherwise closure would
add it. The target contributes one, a contradiction.
Single conclusions and nonnegative multiplicities are essential;
arbitrary signed balances would not give this implication.

\subsection{One product stream}
For labels at most $m$ and $N$ rule instances use linear weights
$1+jX$, and put
\[
 A(X)=\prod_i(1+c_iX),\qquad
 B(X)=(1+X)\prod_i(1+p_iX)(1+q_iX).
\]
Their coefficient sums are bounded by $(m+1)^N$ and
$2(m+1)^{2N}$. Since $m+1\le2^m$ for all $m\in\N$, choose
\[
 M_0=2^{2mN+1},\qquad M=M_0+1.
\]
Every coefficient of $A-B$ has absolute value less than $M$.
Thus equality at $M$ implies polynomial equality: a least nonzero
coefficient would be divisible by $M$ but smaller in absolute value.
Multiplicities of the distinct rational roots $-1/j$ recover the balance.

The recurrence
\[
 S_{i+1}(1+c_iM)=S_i(1+p_iM)(1+q_iM),\qquad
 S_0=M+1,\quad S_N=1
\]
encodes that equality. A proof tree supplies integral accumulator values
from its open obligations. Conversely all factors are positive, so
telescoping is legitimate. The label bound $m$ is chosen before $M$ and
the accumulator. A separate CRT bound on witnesses is introduced later,
preventing a circular bound.

\subsection{Closing the arithmetic}
The predecessor's beta-stream, bounded-CRT certificate, Pell power graph,
factorial graph, and binomial graph replace all the displayed variable
exponents and bounded quantifiers by finite existential polynomial
certificates, with boundary cases included. The normalization changes only
the redundant truncated-subtraction fragment proved above.

Retaining nonlinear wires and substituting nonnegative affine wires gives
the literal quadratic residual system. Its squared sum is quartic.
Only semantic inputs need three-square sign guards: all other wire signs
follow by induction through their defining equations.

\begin{table}[ht]
\centering\small
\begin{tabular}{@{}lrr@{}}
\toprule
Complete direct construction & Incoming & New\\
\midrule
Natural variables &16,036&14,956\\
Unrestricted integer variables &27,916&26,188\\
Degree &4&4\\
Semantic natural inputs &3,960&3,744\\
Retained nonlinear wires &12,076&11,212\\
Quadratic residuals &12,366&11,502\\
Expanded natural monomials &700,865&670,455\\
Expanded integer monomials &758,275&724,768\\
Maximum absolute integer coefficient &320,307,372,009,918&
320,307,372,009,864\\
\bottomrule
\end{tabular}
\caption{A complete global reduction, not only a local witness count.
The integer variable count falls by $1{,}728$, about $6.19\%$.}
\end{table}

\section{Verification and practical use}

\subsection{Exact checks and their limits}
The principal exact checks establish:
\begin{itemize}
\item equality of residual formulas and literal circuits, variable closure,
topological order, liveness, and operation counts;
\item norm identities, denominator clearings, and witness-restoration
identities;
\item exact upper degrees, cancellations, and nonzero highest-degree
coefficients on the actual fixed tuple;
\item complete coefficient byte bindings and retained source identities;
\item every transformed parent arithmetic row and the sign rescaling
identity of the twenty-integer equation;
\item the direct compiler's single changed fragment, all $108$ normalized
power interfaces, all arithmetic compositions, all twelve Horn branches,
global substitutions, and complete final sparse polynomial expansions.
\end{itemize}
The integer direct polynomial has $724{,}768$ nonzero monomials. Its
JSON term format is an integer coefficient and a list of zero-based variable
indices; repeated indices are powers. This is an exact sparse polynomial,
not a symbolic instruction to discover coefficients.

The small equations use ordered registers and declared fixed leaves.
Their circuit files permit only binary addition, subtraction, and
multiplication. Neither representation contains a variable-exponent
operation or oracle.

Exact finite algebra does not establish consistency, unbounded Pell
classification, native computation calibration, or an encoded internal-ZFC
derivation. Those are separate mathematical claims, with the universal
new proofs and inherited dependencies stated explicitly.

\subsection{Reproduction}
Use the archive README and the root verification driver.
The new small-equation checks and direct replay use Python's standard
library. PDFLaTeX rebuilds this report without shell escape or a network.
The final manifest authenticates release files and excludes itself.

The artifact count and the finite degree certificates are reproducible.
PDF bytes can change between TeX distributions without changing the
mathematical source. All frozen coefficient payloads and polynomial
definitions must remain identical.

\subsection{Research implications}
The useful general lesson is that a prescribed residue class can replace
some large radical-separation bounds. The three-root argument does not
automatically extend to four roots: an irrational pair can cancel while
two rational terms remain, creating new residue cases. Writing a larger norm is therefore
not yet a fourteen-variable construction.

The direct branch has a different bottleneck: repeated arithmetic
representations, rather than the fixed logical calculus. Further local
graph improvements should be judged only after composition into the
complete finite equation. A local reduction in variables or degree is not
by itself a global record.

No optimality lower bound or worldwide priority claim is supplied.
The concrete contributions are new calibrated finite equations, explicit
two-way witness proofs, and audited improvements against named inputs.

\appendix
\section{Compact specification of the final fifteen-variable equation}
This is the recommended centered modulus-four equation, of exact degree
$306$ and evaluator size $139$. All inputs are the fifteen integers of
Section~\ref{sec:fifteen}; all other symbols are polynomial registers.
\begin{align*}
 J_0&=3J+1,&q&=\beta J_0+1,&W&=4(3w+1),\\
 X&=qW,&Y&=(2v+1)q^3,&E&=XY,\\
 a&=(X+1)Y,&H&=4a+3,&\Delta&=a^2+H,\\
 F&=(K_0+W)C-(q-1)z,\\
 R&=(q-F)(q^3-q)-C(q^2-1)+(M_C+qM_F)J_0,\\
 k&=R+1+hE,&\eta&=c-Yk,&\zeta&=k-\eta,\\
 m&=XY^2,&B_0&=mk,&S&=k(2m+1),\\
 V&=c(Tf-1)-Rf^2,&s&=X+\gamma H,\\
 \mathcal M&=2acs+s^2-Hc^2-1,\\
 d_{\mathrm g}&=A_{\mathrm c}B_{\mathrm c}FC
                  -(q-A_{\mathrm c}F-B_{\mathrm c}C)^2,\\
 \Pos&=1+p_1^2+p_2^2+p_3(p_3+1),\\
 \mathcal G&=d_{\mathrm g}\eta\zeta
 -(t^2-d_{\mathrm g}-\eta\zeta+1)^2-\Pos,\\
 A_0&=1+m(m+1)k^2,&d&=c^4,&Q&=\Delta(f^2-1),\\
 n&=4(n_0-hYR)+2,&\delta&=Q-1,&D&=d\delta,\\
 L&=n+B_0+4SV,\\
 \overline A&=A_0D,&\overline B&=Q\delta,&
 \overline C&=16S^2d(QV^2-1),\\
 U&=DL^2-\overline A-\overline B-\overline C,\\
 T_3&=U^2-4(\overline A\overline B+
             \overline A\overline C+\overline B\overline C),\\
 N_3&=T_3^2-64\overline A\overline B\overline C\,DL^2,\\
 G&=\mathcal M^2+\mathcal G^2,&P_{15}&=G+N_3+GN_3.
\end{align*}
The uncentered degree-$310$ version changes only $n$ to $4n_0+2$.
The literal counted circuit shares products and uses cancellation
identities; counting the typeset display by eye need not give the same
optimized operation count.

\section{Provenance and reading order}
The snapshot ledger and source notices are in \file{provenance/}.
The small-coefficient direct equation and six-port native equations are
separate components. To assess the main new result, read the root-merger
theorems and parent restorations. For the final consistency interpretation,
follow the retained parent and ordinary calibration proofs. For the
literal finite claims, run the exact checkers.

The ordinary proofs were reviewed within this AI-assisted collaboration,
not by external human referees. The package is a new research contribution;
it does not modify or replace earlier repository packages.

\begin{thebibliography}{12}\small
\bibitem{sat}
ProveIt incoming contribution,
\emph{Saturated Pell certificates and quotient recovery for Con(ZFC)},
integer component, Sections 3--8.
\url{https://github.com/VladimirReshetnikov/ProveIt/blob/43ac596e36c105a0e8b07c2d1f5ae200cae9c056/docs/incoming/ConZFC_Saturated_Certificates_2026-10-24.zip}.
Relevant proofs, equations, coefficient recipe, and licenses are retained.

\bibitem{shears}
ProveIt incoming contribution,
\emph{Polynomial shears and sign recovery for Diophantine consistency}.
\url{https://github.com/VladimirReshetnikov/ProveIt/blob/43ac596e36c105a0e8b07c2d1f5ae200cae9c056/docs/incoming/ConZFC_Polynomial_Shears_2026-10-23.zip}.

\bibitem{projection}
ProveIt incoming contribution,
\emph{Smaller Diophantine consistency equations}.
\url{https://github.com/VladimirReshetnikov/ProveIt/blob/43ac596e36c105a0e8b07c2d1f5ae200cae9c056/docs/incoming/con-zfc-saturation-and-projection-2026-10-23.zip}.

\bibitem{round12}
ProveIt incoming contribution,
\emph{Direct and centered Diophantine consistency equations},
direct fixed-base construction and independent semantic audit.
\url{https://github.com/VladimirReshetnikov/ProveIt/blob/43ac596e36c105a0e8b07c2d1f5ae200cae9c056/docs/incoming/ConZFC_Direct_Centered_Round12_2026-10-24.zip}.

\bibitem{joint22}
ProveIt active research, \emph{Complete centering and degree certificate},
joint22-centered candidate.
\url{https://github.com/VladimirReshetnikov/ProveIt/blob/43ac596e36c105a0e8b07c2d1f5ae200cae9c056/Computability/DiophantineConsistency/candidates/joint22-centered/PROOF.md}.

\bibitem{review}
ProveIt active research, \emph{Wave11 review and size-reduction decisions}.
\url{https://github.com/VladimirReshetnikov/ProveIt/blob/43ac596e36c105a0e8b07c2d1f5ae200cae9c056/Computability/DiophantineConsistency/PROPOSAL_REVIEW_WAVE11.md}.

\bibitem{status}
ProveIt active research, \emph{Native arithmetic and the remaining
calibration contract}; \emph{Object ZFC}.
\url{https://github.com/VladimirReshetnikov/ProveIt/blob/43ac596e36c105a0e8b07c2d1f5ae200cae9c056/Computability/DiophantineConsistency/NATIVE_CALIBRATION52.md}.
Status snapshots are retained in the source ledger.

\bibitem{sun}
Zhi-Wei Sun, \emph{Ten unknowns for Hilbert's tenth problem over the integers},
arXiv:2609.01594, 2026. \url{https://arxiv.org/abs/2609.01594}.

\bibitem{jones}
James P. Jones, \emph{Universal Diophantine equation},
Journal of Symbolic Logic 47 (1982), 549--571.
\url{https://www.cambridge.org/core/journals/journal-of-symbolic-logic/article/113EBF3D928DFC6011C52BECAE81A087}.

\bibitem{carneiro}
Mario Carneiro, \emph{A Lean formalization of Matiyasevi\v{c}'s theorem},
arXiv:1802.01795, 2018. \url{https://arxiv.org/abs/1802.01795}.
Background on formal Pell and exponentiation representations;
the direct compiler uses the supplied pinned sources.
\end{thebibliography}
\end{document}
